\begin{codebox}
\codeline{\textcolor{cmtgray}{\#\ GraphRAG：知识图谱增强的检索生成}}
\codeline{\textcolor{cmtgray}{\#\ GraphRAG:\ Knowledge-Graph-Augmented\ Retrieval\ Generation}}
\codeline{\textcolor{cmtgray}{\#}}
\codeline{\textcolor{cmtgray}{\#\ 朴素RAG检索"文本块"，GraphRAG检索"结构化子图"}}
\codeblank{}
\codeline{\textcolor{cmtgray}{\#\ ==============================}}
\codeline{\textcolor{cmtgray}{\#\ 1.\ 朴素RAG的两个局限}}
\codeline{\textcolor{cmtgray}{\#\ ==============================}}
\codeline{\textcolor{cmtgray}{\#\ 局限1：结构丢失}}
\codeline{\textcolor{cmtgray}{\#\ \ \ 文档切块后，"设备-工序-产品"的关联被切断，}}
\codeline{\textcolor{cmtgray}{\#\ \ \ 检索到的块可能只含关系的一半}}
\codeline{\textcolor{cmtgray}{\#\ 局限2：多跳失效}}
\codeline{\textcolor{cmtgray}{\#\ \ \ 问"P001延期会影响哪些客户订单？"}}
\codeline{\textcolor{cmtgray}{\#\ \ \ 需要\ 产品\(\rightarrow\)工单\(\rightarrow\)订单\(\rightarrow\)客户\ 三跳关联，}}
\codeline{\textcolor{cmtgray}{\#\ \ \ 没有任何单一文本块同时包含整条链路}}
\codeblank{}
\codeline{\textcolor{cmtgray}{\#\ ==============================}}
\codeline{\textcolor{cmtgray}{\#\ 2.\ GraphRAG\ 流程}}
\codeline{\textcolor{cmtgray}{\#\ ==============================}}
\codeline{\textcolor{cmtgray}{\#\ \ \ 用户问题}}
\codeline{\textcolor{cmtgray}{\#\ \ \ \ \ │}}
\codeline{\textcolor{cmtgray}{\#\ \ \ (1)\ 实体识别与链接}}
\codeline{\textcolor{cmtgray}{\#\ \ \ \ \ │\ \ \ \ \ "P001延期影响哪些订单"\ \(\rightarrow\)\ 锚点实体\ mfg:P001}}
\codeline{\textcolor{cmtgray}{\#\ \ \ \ \ │}}
\codeline{\textcolor{cmtgray}{\#\ \ \ (2)\ 子图扩展（k-hop\ /\ 按关系类型走）}}
\codeline{\textcolor{cmtgray}{\#\ \ \ \ \ │\ \ \ \ \ 从P001出发沿\ producedBy/contains/orderedBy\ 扩展2\textasciitilde{}3跳}}
\codeline{\textcolor{cmtgray}{\#\ \ \ \ \ │\ \ \ \ \ 得到：P001\ ←\ W001(工单)\ ←\ O-2026-088(订单)\ ←\ 客户C17}}
\codeline{\textcolor{cmtgray}{\#\ \ \ \ \ │}}
\codeline{\textcolor{cmtgray}{\#\ \ \ (3)\ 子图裁剪与排序}}
\codeline{\textcolor{cmtgray}{\#\ \ \ \ \ │\ \ \ \ \ 按与问题的相关度裁剪（边类型白名单、路径打分），}}
\codeline{\textcolor{cmtgray}{\#\ \ \ \ \ │\ \ \ \ \ 控制在LLM上下文预算内}}
\codeline{\textcolor{cmtgray}{\#\ \ \ \ \ │}}
\codeline{\textcolor{cmtgray}{\#\ \ \ (4)\ 子图序列化\ \(\rightarrow\)\ 注入上下文}}
\codeline{\textcolor{cmtgray}{\#\ \ \ \ \ │}}
\codeline{\textcolor{cmtgray}{\#\ \ \ (5)\ LLM\ 基于子图作答（要求引用图中三元组作为依据）}}
\codeblank{}
\codeline{\textcolor{cmtgray}{\#\ ==============================}}
\codeline{\textcolor{cmtgray}{\#\ 3.\ 子图序列化格式对比}}
\codeline{\textcolor{cmtgray}{\#\ ==============================}}
\codeline{\textcolor{cmtgray}{\#\ 格式A：三元组列表（最忠实，token较省）}}
\codeline{\textcolor{cmtgray}{\#\ \ \ (P001,\ producedBy,\ W001)}}
\codeline{\textcolor{cmtgray}{\#\ \ \ (W001,\ partOf,\ O-2026-088)}}
\codeline{\textcolor{cmtgray}{\#\ \ \ (O-2026-088,\ orderedBy,\ Customer\_C17)}}
\codeline{\textcolor{cmtgray}{\#\ \ \ (O-2026-088,\ dueDate,\ "2026-06-20")}}
\codeblank{}
\codeline{\textcolor{cmtgray}{\#\ 格式B：自然语言化（LLM理解最稳）}}
\codeline{\textcolor{cmtgray}{\#\ \ \ "产品P001由工单W001生产；W001属于订单O-2026-088；}}
\codeline{\textcolor{cmtgray}{\#\ \ \ \ 该订单的客户是C17，交期2026-06-20。"}}
\codeblank{}
\codeline{\textcolor{cmtgray}{\#\ 格式C：JSON（适合函数调用/结构化输出链路）}}
\codeline{\textcolor{cmtgray}{\#\ \ \ \{"entity":"P001","produced\_by":\{"workorder":"W001",}}
\codeline{\textcolor{cmtgray}{\#\ \ \ \ "order":\{"id":"O-2026-088","customer":"C17","due":"2026-06-20"\}\}\}}}
\codeblank{}
\codeline{\textcolor{cmtgray}{\#\ 实践：格式A+B混合——三元组保证可校验，旁附一句话说明降低误读}}
\codeblank{}
\codeline{\textcolor{cmtgray}{\#\ ==============================}}
\codeline{\textcolor{cmtgray}{\#\ 4.\ 检索策略：向量\ vs\ 图\ vs\ 混合}}
\codeline{\textcolor{cmtgray}{\#\ ==============================}}
\codeline{\textcolor{cmtgray}{\#\ 向量检索：召回与问题语义相近的"描述性知识"（维修案例文本）}}
\codeline{\textcolor{cmtgray}{\#\ 图检索：\ \ 精确获取"结构性知识"（关联、路径、聚合）}}
\codeline{\textcolor{cmtgray}{\#\ 混合（推荐）：}}
\codeline{\textcolor{cmtgray}{\#\ \ \ 实体锚定\ \(\rightarrow\)\ 图扩展拿结构\ \(\rightarrow\)\ 对图中实体挂载的文档片段再做向量检索}}
\codeline{\textcolor{cmtgray}{\#\ \ \ （例：拿到BearingWear节点后，再向量检索关联的维修工单文本）}}
\codeblank{}
\codeline{\textcolor{cmtgray}{\#\ 变体：层级社区摘要（Microsoft\ GraphRAG,\ 2024）}}
\codeline{\textcolor{cmtgray}{\#\ \ \ 离线对图做社区检测，逐层生成社区摘要；}}
\codeline{\textcolor{cmtgray}{\#\ \ \ 全局性问题（"本月质量问题集中在哪类设备？"）检索社区摘要而非节点}}
\codeblank{}
\codeline{\textcolor{cmtgray}{\#\ ==============================}}
\codeline{\textcolor{cmtgray}{\#\ 5.\ 与本体的关系：GraphRAG\ 因本体而可靠}}
\codeline{\textcolor{cmtgray}{\#\ ==============================}}
\codeline{\textcolor{cmtgray}{\#\ 没有本体的图谱：边类型随意、方向混乱\ \(\rightarrow\)\ 扩展结果噪声大}}
\codeline{\textcolor{cmtgray}{\#\ 有本体约束的图谱：}}
\codeline{\textcolor{cmtgray}{\#\ \ \ 关系白名单来自本体属性定义（扩展时只走声明过的属性）}}
\codeline{\textcolor{cmtgray}{\#\ \ \ 类型过滤用类层次（"影响订单"只保留\ Order\ 类节点）}}
\codeline{\textcolor{cmtgray}{\#\ \ \ 答案可再走一遍第三层校验（见\ hallucination-control.txt）}}
\codeblank{}
\codeline{\textcolor{cmtgray}{\#\ ==============================}}
\codeline{\textcolor{cmtgray}{\#\ 6.\ 评测要点}}
\codeline{\textcolor{cmtgray}{\#\ ==============================}}
\codeline{\textcolor{cmtgray}{\#\ 多跳问题集：人工构造50\textasciitilde{}100个2\textasciitilde{}3跳问题}}
\codeline{\textcolor{cmtgray}{\#\ \ \ 对比：朴素RAG\ vs\ GraphRAG\ 的答案完整率与事实准确率}}
\codeline{\textcolor{cmtgray}{\#\ 典型结果形态：单跳问题两者接近；2跳以上GraphRAG优势显著}}
\codeline{\textcolor{cmtgray}{\#\ 成本注意：子图扩展的图查询延迟\ +\ 序列化token开销，}}
\codeline{\textcolor{cmtgray}{\#\ \ \ 用边类型白名单与跳数上限控制（一般\ k\ \(\le\)\ 3）}}
\end{codebox}
